\section{An exact selective iteration with a preserved cache}
\label{sec:solver-cache}

The posterior certificate becomes an algorithm when it is paired with a gain
construction that is guaranteed to reach a passing certificate. The canonical
construction in Core 1.0.0 performs exact rational selected-packet updates. It
retains enough residual information to choose the next update without rebuilding
the complete residual after every step. Its final certificate and parameter
representation use outward-rounded dyadic arithmetic. These are separate stages
with separate correctness arguments.

Throughout this section, the state dimension and observation dimension are finite.
The supplied rational model has
\[
 P=LL^{\mathsf T},\qquad R=NN^{\mathsf T}+\sigma I,
 \qquad \sigma>0,
\]
and observation operator $H$. Put
\[
 S=HPH^{\mathsf T}+R,\quad v=y-Hm,\quad
 F(B)=PH^{\mathsf T}-BS,\quad E(X)=\Frob{X}^{2}.
\]
The output packet family $(Q_a)_{a\in A}$ consists of rational self-adjoint,
mutually orthogonal projectors with sum $I$. The input family $(D_b)_{b\in D}$
has the same properties on the observation space. The output family controls
selection and cache maintenance. The input family participates in the independent
final certificate. No commutation relation between $S$ and the input projectors
is assumed.

Preprocessing forms three model-dependent constants:
\begin{equation}
 \tau=\frac{1}{\sigma+\tr S},\qquad T=I-\tau S,
 \qquad w=\frac{\sigma+\lVert v\rVert^{2}}{\sigma^{2}}.
 \label{eq:solver-kernel}
\end{equation}
All are rational. They are the fields of the record
\path{Komorebi.CachedPosterior.Kernel}, produced by
\path{Komorebi.CachedPosterior.prepare}. The correctness theorems use this actual
preparation, rather than an arbitrary record whose fields happen to have the same
names. A wrong transition matrix or weight would invalidate the correspondence
between the stored score and the posterior certificate.

\begin{definition}[Cache validity]
A cache state is a tuple $c=(B,(G_a),(e_a),t)$ satisfying
\begin{equation}
 G_a=Q_aF(B),\qquad e_a=E(G_a),\qquad t=\sum_{a\in A}e_a.
 \label{eq:solver-invariant}
\end{equation}
The cached score is $s(c)=tw$.
\end{definition}

Validity is an exact equality invariant. It does not mean that a stored residual
is approximately near a true residual, nor that its energy is merely an upper
bound. Initialization evaluates $F(B_0)$ once, projects it onto every output
packet, computes every energy, and sums them. The theorem
\path{initialState_valid} follows directly from the definitions. It requires no
prediction about future iterates.

Orthogonal complete packets account for the whole energy:
\[
 \sum_{a\in A}E(Q_aF)=E(F).
\]
Consequently a valid cache has $t=E(F(B))$ and
\begin{equation}
 s(c)=E(F(B))\frac{\sigma+\lVert v\rVert^2}{\sigma^2}=a(B).
 \label{eq:solver-score}
\end{equation}
The right-hand side is the exact residual score used by the posterior error
certificate. The declarations \path{total_eq_residual_energy} and
\path{score_eq} establish these two identifications. The total is therefore a
meaningful global stopping signal even though only one of its summands changes
in an iteration.

For a selected label $a$, the executed update is
\begin{align}
 B'&=B+\tau G_a, &G'_a&=G_aT, &e'_a&=E(G'_a),\nonumber\\
 G'_b&=G_b\quad(b\ne a), &e'_b&=e_b\quad(b\ne a),
 &t'&=t-e_a+e'_a.
 \label{eq:solver-cache-step}
\end{align}
This is the definition \path{Komorebi.CachedPosterior.step}. It computes the new
selected residual from the old selected residual and the precomputed transition.
It retains unselected entries and updates the scalar total by a difference. The
implementation does not define its new state by recomputing $F(B')$ and projecting
that matrix onto every packet.

\begin{proposition}[Preservation by the executed update]
If the output packets satisfy the stated projector assumptions, $c$ is valid,
and the kernel is \path{prepare M}, then the state in
\eqref{eq:solver-cache-step} is valid.
\end{proposition}
\begin{proof}
By cache validity, $G_a=Q_aF(B)$, so the gain update is the exact packet step
$B'=B+\tau Q_aF(B)$. Expand its residual:
\[
 F(B')=F(B)-\tau Q_aF(B)S.
\]
Multiplication on the left by the selected projector gives
\[
 Q_aF(B')=Q_aF(B)-\tau Q_a^2F(B)S
          =G_a(I-\tau S)=G_aT.
\]
For $b\ne a$, mutual orthogonality gives $Q_bQ_a=0$ and hence
$Q_bF(B')=Q_bF(B)=G_b$. This proves every residual equality in the new cache.
Its energy equalities hold because the selected energy is evaluated from its new
residual and the remaining residual-energy pairs are retained. Finally, replacing
one term in a finite sum changes that sum by exactly $-e_a+e'_a$. Thus
$t'=\sum_b e'_b$. All three clauses of validity are preserved.
\end{proof}

This proof is realized by \path{step_valid}, using the algebraic identities
\path{selected_residual} and \path{unselected_residual}. The declarations
\path{step_other_residual} and \path{step_other_energy} expose the unchanged
entries independently. The distinction between proving an invariant for an
executed update and postulating it for a desired state is decisive: the later
termination theorem uses the invariant at every state visited by the program.

The locality is on the output side. For an input projector $D_b$,
\[
 F(B')D_b=F(B)D_b-\tau Q_aF(B)SD_b.
\]
The factor $SD_b$ can mix input directions. The source declaration
\path{input_packet_effect} makes this effect explicit. Complete input packets
can still account for certificate energy, but an input packet decomposition does
not by itself make the iteration independently local in both indices.

\section{Selection, exact trajectory equality, and finite fuel}
\label{sec:solver-termination}

The algorithm chooses a largest stored output energy. Its finite scan has a seed
label and an ordered list. It replaces the incumbent label only when a scanned
label has strictly greater energy. On an equality it preserves the incumbent.
This tie rule appears in the definition \path{winner}, whose comparison is
\path{e a < e b}. The scan may contain repeated labels. Completeness requires
that every output label occurs in the scan, while the seed supplies an inhabitant
of the label type. The proof counts the cardinality of the finite label type,
not the length of a possibly repeated scan.

\path{winner_ge_initial} and \path{winner_ge_member} prove that the final winner
has energy at least that of the seed and of every scanned label. Under complete
coverage, \path{selectPacket_max} therefore supplies a genuine maximum over all
output packet energies. The argument does not prescribe a fixed schedule. The
maximum is selected again from the current state after every update.

\begin{proposition}[Exact selection and gain equality]
For a valid cache, cached selection returns the same label as a scan of the
recomputed energies $E(Q_aF(B))$. Starting from an initialized cache and using the
prepared kernel, its gain after $n$ updates equals the recomputed greedy gain after
$n$ updates, for every $n\in\N$.
\end{proposition}
\begin{proof}
Each stored energy equals the corresponding recomputed energy by
\eqref{eq:solver-invariant}. Both scans therefore receive the same energy
function, the same seed, the same list, and the same strict comparison. Their
outputs are equal, including ties. This is \path{chosen_eq}. The cache gain update
then equals the recomputed greedy step by \path{advance_gain}.

The initial gains agree at $n=0$. Suppose they agree at $n$. The cache at that
index is valid by repeated application of \path{step_valid}. Equality of
selection and equality of the gain step give equality at $n+1$. Induction proves
the entire trajectory identity. The formal statements are
\path{trajectory_valid} and \path{trajectory_gain}.
\end{proof}

This is stronger than convergence to the same limit. At a fixed index, the two
rational gain matrices are identical. As a result, the convergence proof for the
recomputed greedy algorithm transfers directly to the cached execution.

Write $A_{\mathrm{count}}=|A|$ for the positive number of output labels, and set
\begin{equation}
 \rho=1-\frac{\sigma}{\sigma+\tr S},\qquad
 q=\rho^2,\qquad r=1-\frac{1-q}{A_{\mathrm{count}}}.
 \label{eq:solver-rates}
\end{equation}
The model factorization and positive noise floor imply that $S$ is positive
semidefinite with quadratic lower bound $\sigma$. Its trace is nonnegative, and
$\sigma+\tr S$ is a valid positive spectral cap. Thus $0\le\rho<1$, $0\le q<1$,
and $0\le r<1$. No eigenvalue oracle is needed to construct these rational
constants.

The transition contracts every residual matrix:
\[
 E(XT)\le qE(X).
\]
If label $a$ is updated, its residual energy contracts by this inequality and
all other output residuals remain unchanged. Orthogonal energy accounting gives
\[
 E(F(B'))\le E(F(B))-(1-q)e_a.
\]
Because $a$ is a maximum, $e_a\ge E(F(B))/A_{\mathrm{count}}$. It follows that
\begin{equation}
 E(F(B'))\le rE(F(B)),\qquad
 a(B_n)\le r^na(B_0).
 \label{eq:solver-geometric}
\end{equation}
The energy share is \path{selected_energy_share}; one-step contraction is
\path{greedy_energy_contract}; and the score estimate is
\path{greedy_score_geometric}. The estimate allows singular $P$, arbitrary initial
rational $B_0$, rectangular $H$, and rational rotated packets under the same
projector assumptions.

A convergent sequence alone would not define a terminating program with an
explicit natural-number budget. The source supplies a logarithm-free fuel:
\begin{equation}
 N=\left\lceil
 \frac{a(B_0)}{(\varepsilon/2)^2(1-r)}
 \right\rceil_{\N},\qquad \varepsilon>0.
 \label{eq:solver-fuel}
\end{equation}
Here $\lceil\cdot\rceil_{\N}$ is the natural-number ceiling used by
\path{toleranceFuel}. The canonical \path{adaptiveFuel} evaluates this expression
from the actual model, initial gain, tolerance, and number of labels.

\begin{lemma}[Rational fuel sufficiency]
For $0\le r<1$ and $n\ge N$, the geometric envelope obeys
$r^na(B_0)\le\varepsilon^2/4$.
\end{lemma}
\begin{proof}
The source first proves
\[
 r^n\bigl(1+n(1-r)\bigr)\le1.
\]
At the induction step, the inequality
$r(1+(n+1)(1-r))\le1+n(1-r)$ reduces to nonnegativity of
$(n+1)(1-r)^2$. The natural ceiling in \eqref{eq:solver-fuel} gives
$a(B_0)\le n(\varepsilon/2)^2(1-r)$. Multiplying by the nonnegative
$r^n$ and applying the preceding bound yields the required squared tolerance.
The proof remains valid when the contraction factor is zero. The source
statements are \path{geometric_bernoulli}, \path{geometric_fuel_sufficient}, and
\path{greedy_score_at_fuel}.
\end{proof}

This fuel is conservative. The geometric rate is substantially sharper than the
reciprocal envelope used to obtain a simple rational fuel, and the executed search
can stop much earlier. Neither the large natural fuel nor its existence is a
claim about a practical time or memory budget.

\path{search} tests the exact cached score against $\varepsilon^2/4$ at the current
state. If the test fails and fuel remains, it advances once and decrements fuel.
At fuel zero it still tests the last permitted state. This convention permits
success after exactly $N$ updates. \path{search_score} proves that every returned
state passes the score condition. \path{search_origin} proves that a returned
state is an actual trajectory state at some $k\le N$. Finally,
\path{search_complete} proves that a passing score at the last permitted state
forces the search to return some state, possibly earlier.

Combining preserved validity, exact trajectory equality, and fuel sufficiency
therefore proves a finite successful cached search. There is no caller-supplied
claim that an iteration has converged. The computed budget and proved invariant
supply the necessary conclusion internally.

\section{Independent final verification and precision fallback}
\label{sec:solver-verification}

The exact cached score is used to select a gain and establish termination. The
successful gain is still sent to the independent certificate evaluator. That
evaluator reconstructs rounded packet residual obligations from the actual model
and gain. It does not accept the stored cache energies as certificate evidence.
This separation protects successful returns from an incorrect internal stopping
signal. The canonical route checks setup, constructs the outward-rounded
candidate, and executes its tolerance test. Construction and setup imply that
\path{checkCandidate} is true by theorem. A separate untrusted-candidate checker
can recompute exact residual obligations; that is distinct from the canonical
rounded construction route.

The definition \path{solveCached} first checks a positive posterior tolerance,
valid model and packet setup, and complete scan coverage. It prepares the kernel,
initializes the cache, and runs the exact-score search with its computed fuel.
Only the returned gain is passed to \path{certifyGain}. Thus certificate evaluation
is outside the selective update loop. The definition makes one call to this final
verification procedure, although that procedure may try several precisions.

The precision schedule consists of the caller's finite list of probes followed by
one computed fallback. The probes can be empty, repeated, or unordered. Successful
verification records the precision actually used. The appended fallback makes
finite success independent of a lucky externally chosen probe.

Let $h_p=2^{-p}$. The dyadic evaluator propagates rational centers and outward
radii through its arithmetic graph. Its budget calculation supplies a nonnegative
coefficient $C_{\mathrm{ver}}(M,Q,D,B)$ independent of $p$, with
\begin{equation}
 \widehat a_p(B)\le a(B)+C_{\mathrm{ver}}(M,Q,D,B)h_p.
 \label{eq:solver-inflation}
\end{equation}
Here $\widehat a_p$ is the actual rounded candidate score. The full expression for
the coefficient is \path{scoreRoundCoeff}, built from packet residual energy and
innovation-vector magnitude and radius budgets. It bounds the entire score graph,
rather than rounding only its final exact answer. The theorem
\path{rounded_score_inflation} proves \eqref{eq:solver-inflation}.

The common precision selector is the computable integer expression
\begin{equation}
 \operatorname{precisionFor}(C,d)
 =\log^{\N}_2\left(\left\lceil C/d\right\rceil_{\N}\right)+1,
 \qquad d>0,
 \label{eq:solver-precision}
\end{equation}
where $\log^{\N}_2$ denotes the natural-number logarithm \path{Nat.log}, including its
specified value at zero. It is not a floating-point logarithm. For every precision
at least this selected count, \path{coefficient_mesh_le} gives $C2^{-p}\le d$.
The proof uses integer power and ceiling inequalities. The additional bit is a
sufficient choice, not a claim that the bit count is minimal.

For final verification, the solver appends
\[
 p_* = \operatorname{sufficientPrecision}(M,Q,D,B,\varepsilon^2/4).
\]
At its returned search state, $a(B)\le\varepsilon^2/4$. The fallback guarantees
inflation at most another $\varepsilon^2/4$, so
\[
 \widehat a_{p_*}(B)\le\varepsilon^2/2\le\varepsilon^2.
\]
Positive $\varepsilon$ gives a positive inflation allowance. Valid model and
packet setup give the candidate check, while this inequality gives the tolerance
check. The declarations \path{certifyRounded_of_slack} and
\path{certifyGain_complete} establish that the fallback must succeed.

The stricter search threshold supplies room for rounded verification. Searching
only until an exact score happened to equal $\varepsilon^2$ would leave no positive
allowance for an enclosing evaluator. The construction deliberately separates
progress from verification slack. Its inherited geometric fuel is computed for
$\varepsilon/2$, which is exactly the square-root tolerance corresponding to
$\varepsilon^2/4$.

\begin{theorem}[Successful and complete cached gain construction]
For rational factor-coded model data with $\sigma>0$, valid complete rational
output and input packets, a seed and complete output scan, any positive rational
$\varepsilon$, any finite precision probe list, and any rational initial gain,
\path{solveCached} returns a candidate and verification precision. The candidate
and tolerance checks pass, and its gain is an actual greedy iterate at an index
within the computed fuel. Its Joseph Gaussian has Wasserstein distance at most
$\varepsilon$ from the explicit posterior version established by the Bayesian
semantics.
\end{theorem}

The source assembles these conclusions through \path{solveCached_terminates},
\path{cached_solver_gain_origin}, and
\path{cached_solver_reaches_posterior_tolerance}. The semantic posterior used here
is the proved explicit Gaussian kernel version. Its agreement with Bayesian
conditioning is established through disintegration and almost-everywhere agreement
with canonical conditional versions, as discussed in the foundations. The theorem
at the supplied rational observation does not assert uniqueness of every possible
pointwise conditional version at an event of probability zero.

The independent statement \path{solveCached_sound} is logically different from
termination. It says that any successful program return satisfies both Boolean
checks without requiring the caller to supply setup-validity proofs. A returned
candidate has passed the evaluator's checks. Valid input hypotheses are needed
to guarantee that a return exists; they are not external assumptions attached to
the meaning of an already successful return.

\section{Parameter arrays and the meaning of a complete result}
\label{sec:solver-output}

A useful posterior output contains more than a passing gain. For the returned
rational gain $B$, the parameter stage represents
\begin{align}
 \mu_B&=m+B(y-Hm),\nonumber\\
 V_B&=(I-BH)P(I-BH)^{\mathsf T}+BRB^{\mathsf T}.
 \label{eq:solver-joseph-output}
\end{align}
These are the mean and Joseph covariance of the same Gaussian law certified by
the Wasserstein statement. The covariance is positive semidefinite under the model
hypotheses even when the prior is singular. The declaration
\path{approximation_eq_output_gaussian} identifies this parameterized Gaussian
with the approximation law used in the posterior theorem.

The evaluator returns a ball for each mean coordinate and each covariance entry.
For exact rational entry $x$, a ball $(c,r)$ means $|x-c|\le r$. Its radius is an
arithmetic enclosure error for that parameter of the returned approximation.
The output tolerance $\delta>0$ bounds every such radius. It is independent of
the posterior Wasserstein tolerance $\varepsilon$.

Each output arithmetic expression has a precision-independent budget coefficient.
The source sums all mean coefficients and all covariance coefficients to obtain
$C_{\mathrm{out}}$, defined by \path{outputCoefficient}. Every individual coefficient is
nonnegative and no greater than this sum. It chooses
\[
 p_{\mathrm{out}}=\operatorname{precisionFor}(C_{\mathrm{out}},\delta)
\]
and evaluates the output graph at that precision. Since every entry radius is
bounded by its coefficient times $2^{-p_{\mathrm{out}}}$, the same integer
precision bounds all radii by $\delta$. The declarations \path{mean_coeff_le},
\path{covariance_coeff_le}, and \path{evaluate_radii_le} establish this argument.
\path{represent_spec} joins entrywise enclosure and the requested radius bound.

The result record stores verification precision, the checked candidate, parameter
precision, and the evaluated arrays. Its contract \path{ResultValid} contains
seven clauses: the candidate check passes, the tolerance check passes, the
Wasserstein bound holds, the arrays enclose the returned approximation's parameters,
all radii are at most $\delta$, the reported parameter precision equals the actual
precision selector, and the arrays equal the actual evaluator at that precision.
The final two clauses make the reported representation part of the contract,
rather than allowing arbitrary arrays that happen to satisfy a looser description.

\begin{theorem}[Complete first-release output]
Under the valid-input hypotheses above and $\delta>0$,
\path{Komorebi.CachedPosterior.runCached} returns a result satisfying
\path{ResultValid}. Every successful return is sound even without setup-validity
proofs supplied by the caller. A nonpositive output-radius tolerance is rejected.
\end{theorem}

These are precisely \path{runCached_terminates}, \path{runCached_sound}, and
\path{runCached_rejects_nonpositive_output_tolerance}. The complete contract is
finite-dimensional inference for one observation update. It includes singular
priors, rectangular observation operators, gain origin, computable iteration fuel,
verification precision, and parameter precision.

Two interpretations must remain separate. The intervals enclose the parameters
of the certified Joseph Gaussian. They do not certify another Gaussian obtained
by treating interval centers as exact mean and covariance. In particular, rounded
covariance centers are not independently proved positive semidefinite. Likewise,
$\delta$ is not a coordinate error bound against the exact posterior. The
Wasserstein guarantee compares laws; the interval guarantee encloses the parameters
of one of those laws. A theorem combining either guarantee with a new operation
must account for that operation explicitly.

The program is mathematically finite for every positive rational pair
$(\varepsilon,\delta)$ under its declared assumptions. A reference implementation
can still encounter a time, memory, or user-supplied resource cap. The Python tools
report \path{RESOURCE_LIMIT} when their caps are reached. Such a cap is a limit of
the execution environment, not a counterexample to the finite-success theorem.
The result contract also contains no sampler and no guarantee for an external
floating-point implementation.

\section{A complete coordinate execution}
\label{sec:solver-coordinate}

The coordinate fixture has a two-dimensional state and one-dimensional
observation:
\[
 P=\begin{pmatrix}4&0\\0&1\end{pmatrix},\quad H=\begin{pmatrix}1&1\end{pmatrix},
 \quad R=1,\quad m=\begin{pmatrix}0\\0\end{pmatrix},\quad y=1.
\]
It is represented by $L=\operatorname{diag}(2,1)$, $N=0$, and $\sigma=1$.
The two output projectors select the coordinate axes, and the single input
projector is $1$. Start with $B_0=0$, seed $0$, scan $[0,1]$, probes $[4,12]$,
$\varepsilon=1/20$, and $\delta=1/1000$.

Here $S=6$, $\tau=1/7$, $T=1/7$, and $w=2$. The initial residual is $(4,1)^{\mathsf T}$,
with packet energies $16$ and $1$. Updating one coordinate multiplies its residual
by $1/7$ and its energy by $1/49$. The other coordinate remains exactly unchanged.
The maximum-energy schedule therefore alternates between labels $0$ and $1$.
The complete pre-update trace is
\begin{center}
\begin{tabular}{rrrrr}
\toprule
Updates already made & Selected label & $e_0$ & $e_1$ & $t$\\
\midrule
0 & 0 & $16$ & $1$ & $17$\\
1 & 1 & $16/49$ & $1$ & $65/49$\\
2 & 0 & $16/49$ & $1/49$ & $17/49$\\
3 & 1 & $16/2401$ & $1/49$ & $65/2401$\\
4 & 0 & $16/2401$ & $1/2401$ & $17/2401$\\
5 & 1 & $16/117649$ & $1/2401$ & $65/117649$\\
\bottomrule
\end{tabular}
\end{center}

After six updates, each coordinate has been selected three times. The geometric
sum in its gain gives
\[
 B_6=\begin{pmatrix}228/343\\57/343\end{pmatrix},\qquad
 F(B_6)=\begin{pmatrix}4/343\\1/343\end{pmatrix},\qquad
 a(B_6)=\frac{34}{117649}.
\]
The strict search threshold is $\varepsilon^2/4=1/1600$. At the previous state,
$a(B_5)=130/117649>1/1600$, while $a(B_6)<1/1600$. Thus six is the actual update
count for this stopping rule. The formal fuel is $410093$; the program uses only
six of these allowed updates. This large difference illustrates the conservative
fuel construction, not an ambiguity in the observed trace.

At verification precision $4$, the rounded score is $25/16$ and the attempt fails.
At precision $12$, it is $481/1048576$, which is below the required final threshold
$\varepsilon^2=1/400$. The solver therefore returns verification precision $12$.
A failed coarse enclosure is a failed certificate attempt, not evidence that the
gain itself is inaccurate. Sufficiently high precision bounds the enclosure inflation.

The exact posterior mean is $(2/3,1/6)^{\mathsf T}$, while the returned Joseph mean
is the gain $B_6$ because $m=0$ and $y=1$. Parameter output uses precision $32$.
Its mean centers are
\[
 \begin{pmatrix}2854963683/4294967296\\89217615/536870912\end{pmatrix}
\]
with radius $1/4294967296$ in each coordinate. The covariance centers are
\[
 \begin{pmatrix}
 5726720411/4294967296 & -1431643597/2147483648\\
 -1431643597/2147483648 & 3579145495/4294967296
 \end{pmatrix},
\]
and the corresponding radii are
\[
 \begin{pmatrix}
 1/536870912&7/4294967296\\
 7/4294967296&1/536870912
 \end{pmatrix}.
\]
Every radius is smaller than $1/1000$. These fractions describe output arithmetic,
not uncertainty intervals for model estimation. They enclose exactly the Joseph
mean and covariance at $B_6$.

The exact rational matrix enclosed by those covariance balls is
\[
 V_{B_6}=\frac{1}{117649}
 \begin{pmatrix}156868&-78432\\-78432&98041\end{pmatrix}.
\]
It can be checked directly by substituting $B_6$ into
\eqref{eq:solver-joseph-output}. Another useful check starts from the exact
posterior covariance
$V_* = \left(\begin{smallmatrix}4/3&-2/3\\-2/3&5/6\end{smallmatrix}\right)$.
If $B_*=(2/3,1/6)^{\mathsf T}$ is the exact gain, expansion of the Joseph
expression gives
\[
 V_{B_6}=V_*+6(B_6-B_*)(B_6-B_*)^{\mathsf T}.
\]
The gain discrepancy is
$B_*-B_6=(2/1029,1/2058)^{\mathsf T}$, which produces the displayed matrix.
This second calculation checks the sign of the off-diagonal entries and explains
why the approximate covariance has an additional positive semidefinite term.
The posterior accuracy certificate concerns the complete mean-covariance law,
while the dyadic output stage faithfully represents these rational parameters.
The larger output budget can still choose many more bits than the smallest
possible entrywise enclosure would require. Its purpose is a guaranteed common
precision, not an optimal encoding of this particular covariance.

The module \path{GeneratedCachedCoordinate.lean} defines this actual model,
initial cache, six-step cache trajectory, final gain, certificate, and result.
Its declarations check the computed fuel, final gain, exact and rounded scores,
each of the six chosen labels, solver return, parameter precision, every displayed
center and radius, and the full output return. Its theorem \path{output_valid}
applies \path{runCached_sound} to the proved return. This is a checked finite
program example, while the universal termination theorem covers all valid inputs
in the declared class.

An earlier recomputed selective solver verified after each update and could stop
at five updates with a passing rounded certificate. That is a different stopping
policy. The canonical cached solver first waits for the stricter exact-score
margin and then performs final verification. A recomputed baseline using this
same strict margin returns the same six-step gain. Comparisons between cached and
recomputed execution must align their stopping policy before comparing work.

\section{A singular rotated execution}
\label{sec:solver-rotated}

The rotated fixture makes the prior singular and the packets non-coordinate:
\[
 P=J=\begin{pmatrix}1&1\\1&1\end{pmatrix},\quad H=I,
 \quad R=I,\quad m=0,\quad y=\begin{pmatrix}1\\-1\end{pmatrix}.
\]
It uses the rectangular prior factor $L=(1,1)^{\mathsf T}$, zero noise factor,
and $\sigma=1$. Both output and input packet families are
\[
 Q_+=\frac12\begin{pmatrix}1&1\\1&1\end{pmatrix},\qquad
 Q_-=\frac12\begin{pmatrix}1&-1\\-1&1\end{pmatrix}.
\]
These project onto the common and antisymmetric directions. They are rational,
self-adjoint, orthogonal, and complete. The prior has rank one, but the innovation
matrix $S=I+J$ is positive definite because of the positive noise floor.

Take $B_0=0$, seed $0$, scan $[0,1]$, probes $[4,12]$, $\varepsilon=1/4$, and
$\delta=1/1000$. Here $\tr S=4$, $\tau=1/5$, and
$T=I-(I+J)/5$. The initial residual is $J$, entirely in the $Q_+$ output packet.
Since $JS=3J$, each selected update multiplies this residual by $2/5$. The
$Q_-$ residual remains zero. Thus every selected label is $0$ and
\[
 (e_+,e_-)=(4(4/25)^n,0),\qquad B_n=\frac{1-(2/5)^n}{3}J.
\]
The first four pre-update energy pairs are $(4,0)$, $(16/25,0)$,
$(64/625,0)$, and $(256/15625,0)$.

After four updates,
\[
 B_4=\frac{203}{625}J,\qquad F(B_4)=\frac{16}{625}J,
 \qquad a(B_4)=\frac{3072}{390625}.
\]
The score weight is $w=3$, because $\lVert y\rVert^2=2$. The strict margin is
$\varepsilon^2/4=1/64$. At the previous state the score is $768/15625>1/64$;
after the fourth it is below $1/64$. The computed fuel is $4267$. Precision $4$
produces score $435/64$ and fails; precision $12$ produces $9111/1048576$ and
passes the final threshold $1/16$.

The exact posterior mean is zero because $Jy=0$, and the returned approximation's
mean is also zero. Its Joseph covariance is
\[
 V_{B_4}=\frac{130379}{390625}J.
\]
The exact posterior covariance is $J/3$. Both laws remain supported in the prior's
one-dimensional common-direction subspace. This example therefore combines a
singular prior, rotated packet selection, a nontrivial covariance approximation,
and a zero mean without using an inverse of the singular prior covariance.

At output precision $32$, each mean center is zero with radius $1/2147483648$.
Every covariance entry has center $179191533/536870912$ and radius
$7/4294967296$. All output radii satisfy $\delta=1/1000$. The equality of these
centers reflects the exact common-direction structure of the returned Joseph
covariance. It does not turn the rounded center matrix into the object of the
posterior theorem.

\path{GeneratedCachedRotated.lean} checks the same kinds of finite obligations as
the coordinate module: fuel, chosen labels, gain, scores, actual return, parameter
precision, arrays, and \path{ResultValid}. The example is particularly useful for
testing whether a presentation has accidentally assumed an invertible prior or
coordinate-only packets. Neither assumption belongs to the solver contract.

\section{The rounded-cache counterexample and the separate rounded solver}
\label{sec:solver-rounding-boundary}

Exact cache validity cannot be reused unchanged after independently rounding its
components. The release contains a scalar counterexample with
$P=H=R=1$, $\sigma=1$, $m=0$, $y=1$, and $B_0=0$. Then $S=2$, $\tau=1/3$,
$w=2$, and the first exact update is
\[
 B_1=1/3,\qquad F(B_1)=1/3.
\]
The naive route separately rounds the gain center and residual center downward
to zero fractional bits. Both become zero. It then computes the stored energy
from the rounded residual, obtaining a stored score of zero.

The actual residual at its stored gain is instead
\[
 F(0)=1,
 \qquad a(0)=2.
\]
Thus the state claims exact internal convergence while its true residual is
nonzero. The equality $G=QF(B)$ has been destroyed by rounding independently
related quantities. The source module \path{CachedRoundingBoundary.lean} proves
\path{naive_score_zero}, \path{true_score_two}, and \path{naive_cache_invalid}.
For $\varepsilon=1/10$ and probe $[12]$, it also proves
\path{naive_final_verification_rejected}: \path{certifyGain} returns no candidate.

This result identifies both a failure and a protection. The naive rounded cache
cannot support the exact-cache termination proof. The independent final verifier
still rejects its false stopping signal, preserving the meaning of successful
returns. The example does not rule out a rounded selective cache with propagated
error bounds. Such a construction needs an enclosure relation, an error budget
for each selected update, and a selection argument that tolerates uncertain
energies. Those obligations are a separate extension.

Core 1.0.0 already retains a different, completed rounded solver. Its update is a
genuinely rounded full-gain Richardson step. It neither claims exact selective
cache invariants nor selects a largest cached packet. Its declarations expose
the whole-update defect instead. If $\widetilde B_{n+1}$ is the rounded update
and $B_{n+1}^{\mathrm{ideal}}$ is the exact full step at the current gain, then
\[
 F(\widetilde B_{n+1})=F(\widetilde B_n)T
 -(\widetilde B_{n+1}-B_{n+1}^{\mathrm{ideal}})S.
\]
This is \path{rounded_residual_step}. The defect is explicitly included in the
tracking and termination argument.

Writing $d$ for the number of gain entries and $h_t=2^{-t}$, the source establishes
an all-iteration tracking bound
\[
 E(\widetilde B_n-B_n^{\mathrm{exact}})
 \le\frac{6d h_t^2}{(1-q)^2}.
\]
Its residual-score comparison has the form
\[
 a(\widetilde B_n)\le2a(B_n^{\mathrm{exact}})+C_{\mathrm{round}}h_t^2,
 \qquad
 C_{\mathrm{round}}=
 \frac{12dE(S)w}{(1-q)^2}.
\]
The source names are \path{roundedGain_tracking}, \path{roundoffScoreCoeff}, and
\path{rounded_score_compare}. The first term is ideal progress. The second is a
precision-dependent floor. Both are needed.

The chosen update precision \path{updateAccuracy} places the floor below
$\varepsilon^2/8$. The rounded fuel \path{roundedFuel} drives the exact full-step
score below $(\varepsilon/4)^2=\varepsilon^2/16$. The factor of two in the
comparison therefore gives another $\varepsilon^2/8$, so the actual rounded
iterate reaches score at most $\varepsilon^2/4$. The same independent final
precision fallback can certify it. \path{solveRounded_terminates} proves finite
successful return, and \path{rounded_solver_gain_origin} identifies the result
as an actual rounded iterate within its declared fuel.

There are consequently two completed branches: exact selective updates with an
exact cache, and genuinely rounded full-gain updates with an explicit perturbation
bound. A rounded selective cache combining their advantages remains an extension.
No proof for one branch is presented as a theorem about the other.

\section{Measured work, execution boundaries, and extensions}
\label{sec:solver-performance}

The cached representation reduces the work needed to maintain packet residuals
relative to recomputing every packet at every selected update. It stores the
projected residual matrices and energies, retains unselected objects, and updates
the total by one scalar difference. Selection still scans the stored energies.
The amount of storage depends on the number and shape of the represented packet
matrices. Rational numerator and denominator growth also matters in execution.
The algebraic locality theorem alone is not a formal bit-complexity estimate.

The archived bounded benchmark compares cached selective, recomputed selective,
and full Richardson reference implementations on seven cases of dimensions at
most eight. All use exact fraction arithmetic, the same zero-skipping matrix
backend, the same strict exact-score stopping margin, and common requested
posterior and output tolerances. All include setup, final rounded certification,
and parameter output. The recomputed and cached selective branches return
identical candidates. The full baseline maintains a full residual by recurrence
and avoids unnecessary packet scans.

One warmup precedes five repetitions with rotating method order. Archived median
producer times and selective-iteration ratios are
\begin{center}
\begin{tabular}{lrrrr}
\toprule
Case & Cached ms & Recomputed ms & Full ms & Iteration ratio\\
\midrule
coordinate & 2.38 & 2.46 & 2.33 & $1.94$\\
rotated & 4.38 & 4.45 & 4.28 & $2.16$\\
sparse4 & 13.84 & 14.52 & 13.79 & $2.51$\\
dense4 & 16.34 & 17.96 & 16.34 & $3.11$\\
sparse8 & 171.98 & 174.30 & 169.87 & $4.62$\\
dense8 & 200.69 & 235.15 & 195.44 & $5.61$\\
conditioned4 & 19.16 & 21.36 & 17.46 & $3.08$\\
\bottomrule
\end{tabular}
\end{center}
The final column is recomputed-selective iteration time divided by cached
iteration time. It excludes setup, verification, and output. The producer columns
include these stages and exclude export and Lean checking. These are archived
reference measurements, rather than fresh timing measurements from this
publication session.

The cache improves selective iteration in these cases, yet these measurements
establish no overall advantage over the efficient full baseline. Several cases
require substantially more selective updates than full updates. The dense8 case,
for example, takes $67$ cached selective updates and $9$ full updates. Final
certification and output account for a large part of total production cost. The
bounded result supports a narrower conclusion: preserving a packet cache avoids
unnecessary selective recomputation, while end-to-end performance depends on more
than the selective update loop.

Every method-case pair also exported an explicit final candidate certificate,
for $21$ checked benchmark outputs. Their archived Lean checks certify those
explicit rational gains and posterior tolerances. They do not certify the Python
timing trace, counters, or numeric parameter arrays from the benchmark runs. The
two generated coordinate and rotated modules have the stronger finite-program and
output-array checks described above. A universal theorem about the mathematical
output evaluator is also distinct from a universal refinement theorem for a
Python implementation of that evaluator.

The archived per-instance Lean checking measurements range from seconds to tens
of seconds and dominate millisecond producer times. Each includes a fresh pinned
Lean process, import, elaboration, and kernel checking with common dependencies
already compiled. The dependency build is accounted for separately. One checking
time per export cannot reliably rank small differences. The source reports preserve
these distinctions, rather than combining an iteration ratio with an unrelated
claim of total acceleration.

The producer/exporter boundary is similarly precise. Python parses user data,
executes a reference algorithm, and emits a Lean instance containing explicit
rational data. Successful checking establishes the statement about that emitted
instance. A universal theorem identifying arbitrary parser input with the emitted
model is outside this release. The output record is a mathematical representation
with checked arithmetic semantics, rather than a claim about every external file
format or hardware arithmetic instruction.

The most direct extensions now have concrete obligations. A rounded selective
cache must propagate residual and energy enclosures and prove progress under its
actual selection rule. Approximate maximum selection needs a quantitative energy
share bound. A changing packet hierarchy must preserve global accounting across
refinement and coarsening. Uncertain model inputs need their own error budget in
both the certificate and the posterior semantics. Sequential filtering must
account for repeated updates and transport of earlier approximation error.
Large-scale performance work must measure preprocessing, final certification,
output representation, rational bit growth, memory, and checking cost as well as
iteration kernels.

These extensions do not leave the first-release theorem unfinished. The closed
contract remains exact rational, finite-dimensional, single-update Gaussian
inference with complete fixed packets, positive tolerances, a terminating gain
construction, independent dyadic verification, and enclosed parameters of the
certified Joseph law. The rounded-cache counterexample and bounded performance
result make the limits of that contract inspectable and give specific targets for
the next stage of development.
